\documentclass[aspectratio=169]{beamer}
\usetheme{Madrid}
\usecolortheme{dolphin}
\setbeamertemplate{navigation symbols}{}
\setbeamertemplate{footline}[frame number]
\setbeamersize{text margin left=6mm,text margin right=6mm}
\usepackage[T1]{fontenc}
\usepackage{lmodern}
\usepackage{booktabs,tabularx,array}
\IfFileExists{xurl.sty}{\usepackage{xurl}}{\usepackage{url}}
\hypersetup{hidelinks,pdftitle={Attempt 4 Literature Review - What We Actually Tried and Why It Stopped}}
\newcommand{\paper}[1]{\textbf{#1}}
\title{Attempt 4 Literature Review}
\subtitle{\texorpdfstring{What we actually tried and why it stopped\\[2mm]\small Corrected next-meeting draft --- not yet presented}{What we actually tried and why it stopped --- Corrected next-meeting draft, not yet presented}}
\author{}
\date{2026-07-22}
\begin{document}
\begin{frame}\titlepage\end{frame}

\begin{frame}{1. First, Correct the Main Misunderstanding}
The original Attempt 4 was \textbf{not} a progressive-pruning proposal.
\begin{columns}[T]
\column{0.47\textwidth}
\begin{block}{Original Attempt 4}
Use the full fixed-width word once. Allow arbitrary scalar level counts and encode them with a mixed-radix address.
\end{block}
\column{0.47\textwidth}
\begin{block}{Later prefix successor}
Read a prefix for a coarse pass, then read the full word for a more accurate pass.
\end{block}
\end{columns}
\begin{alertblock}{Correction}These are different research questions with different stop reasons.\end{alertblock}
\end{frame}

\begin{frame}{2. Original Intuition: Use the Word More Fully}
Two scalar coordinates share one four-bit word, so there are 16 possible addresses.
\begin{center}
\begin{tabular}{ccl}
ordinary & $2$ bits $+$ $2$ bits & $4\times4=16$ states\\
ordinary & $1$ bit $+$ $3$ bits & $2\times8=16$ states\\
Attempt 4 & arbitrary $K_1,K_2$ & $K_1K_2\leq16$
\end{tabular}
\end{center}
Ordinary allocation forces each level count to be a power of two. Attempt 4 also permits choices such as $(K_1,K_2)=(3,5)$, which uses 15 addresses.
\begin{block}{Intuition}Three useful levels on one coordinate and five on another may fit the data better than every legal power-of-two allocation.\end{block}
\end{frame}

\begin{frame}{3. Running Example: What $(3,5)$ Proves}
\begin{center}
coordinate 1: $\{-1,0,1\}$ \hspace{8mm} coordinate 2: $\{-2,-1,0,1,2\}$
\end{center}
\begin{tabularx}{\textwidth}{>{\raggedright\arraybackslash}X c c}
\toprule Allocation & Joint states & Error per vector\\
\midrule arbitrary scalar product $(3,5)$ & 15 & $0$\\
best power-of-two product $(4,4)$ & 16 & $0.1$\\
unrestricted 16-center two-dimensional codebook & 16 & $0$\\
\bottomrule
\end{tabularx}
\begin{block}{Claim ceiling}The example proves that the power-of-two restriction can waste choices. It proves no natural-data, Recall, or speed advantage. The unrestricted block codebook also fits exactly.\end{block}
\end{frame}

\begin{frame}[fragile]{4. How the Original Code Would Work}
For labels $z_1,z_2$, store one mixed-radix address:
\begin{verbatim}
u = z1 + K1 * z2
\end{verbatim}
For $(K_1,K_2)=(3,5)$, addresses 0--14 are used and address 15 is unused.
\medskip
For each query, build one group table:
\begin{verbatim}
T[u] = distance(query group, reconstruction u)
\end{verbatim}
The scan reads one fixed word and performs one lookup per group.
\begin{alertblock}{Not part of the mechanism}No prefix, progressive decoding, or first-stage candidate pruning.\end{alertblock}
\end{frame}

\begin{frame}{5. What Prior Work Already Covers}
\paper{Muresan--Effros (2002)}: optimal one-dimensional quantizers with an integer number of levels, including non-power-of-two counts.
\medskip
\paper{Brandt (2010)}: learned scalar quantizers, integer bit allocation, packed database words, query tables, and nearest-neighbor scan.
\medskip
\paper{Finite Scalar Quantization (2024)}: product codebooks built from small scalar level sets, including non-power-of-two factors.
\medskip
\paper{Mixed-radix indexing}: the standard mapping from several finite alphabets to one integer address.
\begin{block}{Novelty boundary}Arbitrary level counts and mixed-radix packing alone are not enough. The remaining claim must be a measured database accuracy--cost advantage.\end{block}
\end{frame}

\begin{frame}{6. BAPQ and Quicker ADC: Close, but Different}
\paper{Adaptive Bit Allocation Product Quantization (BAPQ, 2016)} gives different integer bit counts to independent subspaces. Five bits still means exactly 32 codewords, not an arbitrary count such as 27.
\medskip
\paper{Quicker ADC (2021)} studies four-, five-, and six-bit product subcodes, packed layouts, split lookup tables, and processor cost. Its sizes also remain powers of two.
\begin{block}{Why they matter}They do not implement $(3,5)$ mixed radix, but they are mandatory strong baselines for allocation, packing, lookup, and scan speed.\end{block}
\scriptsize BAPQ = adaptive bit allocation product quantization; ADC = query-to-code distance lookup.
\end{frame}

\begin{frame}{7. Dense Rates and Arbitrary Levels Are Crowded}
\begin{itemize}
  \item \paper{FSQ (2024)} already uses products of small scalar level sets.
  \item \paper{Q-Palette (2025)} studies fractional-rate quantizers and mixed allocation.
  \item \paper{FibQuant (2026)} studies dense fixed-rate random-access vector coding and exposes the advantage of a vector code over a scalar-product restriction for its canonical source.
\end{itemize}
These papers do not collectively contain the exact database scan design.
\begin{alertblock}{Paper-level burden}``More rate choices'' is not a contribution by itself. Attempt 4 must move a measured database-systems trade-off.\end{alertblock}
\scriptsize FSQ = finite scalar quantization.
\end{frame}

\begin{frame}{8. Product Quantization and Block VQ Are Two Controls}
\begin{tabularx}{\textwidth}{>{\raggedright\arraybackslash}p{0.30\textwidth} X}
\toprule Control & Question answered\\
\midrule optimized product quantization & Does an ordinary power-of-two product code already obtain the same accuracy and scan speed?\\
same-capacity two-dimensional block vector quantization & Does the scalar-product restriction lose too much by forbidding general two-dimensional centers?\\
\bottomrule
\end{tabularx}
\medskip The block codebook is more expressive. Attempt 4 needs a measured reason to retain factorization: for example, lower training cost, less model state, or cheaper table construction.
\scriptsize VQ = vector quantization.
\end{frame}

\begin{frame}{9. What We Actually Established}
\begin{description}
  \item[A4-0] The synthetic $(3,5)$ example passed. This validated the instrument only.
  \item[A4-1S] The exact frozen construction stopped on its cost gate before natural-data work.
  \item[V2] No valid terminal scientific record was produced; runtime remains unknown.
  \item[A4-R0] The representation cannot be inserted as an unchanged SAQ encoder swap.
\end{description}
\begin{alertblock}{Still unmeasured}Natural-data prevalence, Recall, throughput, optimized product quantization, and trained block vector quantization.\end{alertblock}
\end{frame}

\begin{frame}{10. What the 24-Hour Stop Means}
\begin{tabular}{lr}
\toprule Recorded work & CPU time\\
\midrule exact scalar construction & $4.64$ hours\\
canonical evidence serialization & $5.03$ hours\\
completed 61-of-128-coordinate prefix & $9.668$ hours\\
registered full-pipeline lower bound & $24.17$ hours\\
\bottomrule
\end{tabular}
\medskip
The frozen ceiling was $24.0$ CPU-hours, so that registered procedure stopped.
\begin{block}{Correct interpretation}This is a cost result for one exact solver-and-evidence pipeline. It is not a lower bound for every possible implementation of arbitrary-cardinality quantization.\end{block}
\end{frame}

\begin{frame}{11. Why It Was Not an Unchanged SAQ Swap}
Attempt 4 needs learned centers, per-group cardinalities, mixed-radix labels, unused-address handling, new distance tables, and a joint-label query consumer.
\medskip
Existing SAQ stores per-coordinate bitplanes and reconstructs its estimate using a uniform step size and a per-vector rescaling factor.
\begin{alertblock}{R0 decision}\texttt{NO\_GO\_UNCHANGED\_SAQ\_COMPATIBILITY}: equal payload bytes do not give the bytes the same meaning. Integration requires a new representation and query path.\end{alertblock}
This does not establish that arbitrary cardinalities are ineffective.
\end{frame}

\begin{frame}{12. The Prefix Proposal Came Later}
Only after R0 did the project define a broader successor:
\begin{center}
\begin{tabular}{c|c}
first $p$ bits & all $B$ bits\\
coarse lookup and filtering & accurate lookup
\end{tabular}
\end{center}
It added a learned label permutation, nested prefix partition, and two query stages. Its protocol explicitly called itself \textbf{a new direction}.
\begin{block}{Why progressive papers appear}They evaluate this later successor, not the original mixed-radix intuition.\end{block}
\end{frame}

\begin{frame}{13. Why Only the Prefix Successor Was Closed}
\begin{itemize}
  \item \paper{Riskin et al. (1994)} add progressive reconstructions after fixing a fine vector quantizer.
  \item \paper{Derived Codebooks (2019)} derive a coarse grouping from a fine product code, then scan and refine.
  \item \paper{Polysemous Codes (2016)} learn labels for cheap filtering while retaining accurate product-code lookup.
\end{itemize}
Combining these mechanisms with a known arbitrary-level product code left no new constraint or algorithm.
\begin{alertblock}{S0 decision}\texttt{NO\_GO\_DIRECT\_COMPOSITION} closes the prefix successor. It is not the stop reason for original Attempt 4.\end{alertblock}
\end{frame}

\begin{frame}{14. Keep the Two Conclusions Separate}
\scriptsize
\begin{tabularx}{\textwidth}{>{\raggedright\arraybackslash}p{0.23\textwidth} >{\raggedright\arraybackslash}p{0.33\textwidth} X}
\toprule Direction & What stopped it & What remains unknown\\
\midrule original arbitrary-cardinality Attempt 4 & frozen exact-pipeline cost; incompatible with unchanged SAQ & natural-data gain, Recall, speed, fair product/block comparison\\
later prefix-code successor & direct composition of known mechanisms & implementation not justified\\
\bottomrule
\end{tabularx}
\normalsize
\begin{block}{Accurate summary}The original coding intuition remains experimentally unresolved. A later prefix reformulation was separately rejected on prior-work grounds.\end{block}
\end{frame}

\begin{frame}{15. If We Reconsider Original Attempt 4}
A fresh decision should return to the original question and require:
\begin{enumerate}
  \item a cheaper construction method with controlled approximation error;
  \item separately optimized power-of-two product and same-capacity block-vector controls;
  \item complete accounting of model bytes, table construction, encoding time, and scan cost; and
  \item a frozen test of the Recall--throughput frontier.
\end{enumerate}
\begin{alertblock}{Current boundary}This is the unresolved evidence list, not execution authorization. It does not reopen the stopped branch.\end{alertblock}
\end{frame}

\begin{frame}[allowframebreaks]{References and Decision Records}
\tiny
\begin{itemize}
  \item Muresan and Effros. ``Quantization as Histogram Segmentation.'' DCC 2002.
  \item Brandt. ``Transform Coding for Fast Approximate Nearest Neighbor Search in High Dimensions.'' CVPR 2010.
  \item Mentzer et al. ``Finite Scalar Quantization.'' ICLR 2024. \url{https://arxiv.org/abs/2309.15505}
  \item Ge et al. ``Optimized Product Quantization.'' CVPR 2013.
  \item Guo et al. ``Adaptive Bit Allocation Product Quantization.'' Neurocomputing, 2016. \url{https://doi.org/10.1016/j.neucom.2015.07.062}
  \item Andre et al. ``Quicker ADC.'' TPAMI 2021. \url{https://arxiv.org/abs/1812.09162}
  \item Lee and Song. ``Q-Palette.'' 2025. \url{https://arxiv.org/abs/2509.20214}
  \item Lee and Kim. ``FibQuant.'' 2026. \url{https://arxiv.org/abs/2605.11478}
  \item Riskin et al. ``Index Assignment for Progressive Transmission of Full Search Vector Quantization.'' IEEE TIP, 1994. \url{https://doi.org/10.1109/83.287025}
  \item Andre et al. ``Derived Codebooks for High-Accuracy Nearest Neighbor Search.'' 2019. \url{https://arxiv.org/abs/1905.06900}
  \item Douze et al. ``Polysemous Codes.'' ECCV 2016. \url{https://arxiv.org/abs/1609.01882}
  \item Original formulation: \texttt{saq\_attempt4\_arbitrary\_cardinality\_related\_work\_and\_gate\_2026\_07\_13.md}.
  \item R0: \texttt{saq-a4-r0-static-gate@617ad25}. Prefix successor S0: \texttt{saq-a4-prefix-novelty-gate@7f4c001}.
\end{itemize}
\end{frame}
\end{document}
